% ============================================================
% Chapter 5: Development
%   Section: News Acquisition
%   Subsection: Trend Signals
% Included from chapters/05-development/02-news-acquisition/news-acquisition.tex
% Source: docs/decisions/trends.md; backend/src/services/trends/;
% backend/data/trends/; the feeds saved on 2026-10-26 in
% backend/tests/fixtures/trends/.
% ============================================================

\subsection{Trend Signals}
\label{sec:dev-trend-signals}

The trending lens (Section~\ref{sec:design-selection}) reads
\emph{signals}: public feeds that say what is spreading right now. A
signal is not a news source. A source is read for its articles and
rated, by domain, as evidence; a signal is never evidence, and Google
Trends is one domain for every country. Signals therefore have their own
folder of declarations, one per signal, and their own reader
(\path{backend/src/services/trends/}).

\paragraph{The signals.} Google Trends' ``Trending now'' feed for Spain
and for the United States gives the country's ten rising searches, each
with its approximate searches and three news articles about it; a trend
is the search, and its articles are Google's. The day's top posts of
r/UpliftingNews give good news that is already spreading: a trend is a
post, and its article is the page it was submitted with, never the
discussion thread. Upworthy and Know Your Meme are read as ordinary feeds
whose order is the signal. The feeds are read concurrently through the
same fetcher as every page (Section~\ref{sec:dev-scraping}): the URL
guard, a timeout, redirects checked hop by hop. Each read is counted on
the scraper's page, and a signal that cannot be read is reported with
the reason, never raised. Google's feed is parsed with the standard
library's XML parser, because the general feed parser keeps only the
last of a repeated namespaced element and would drop two of each trend's
three articles.

\paragraph{A person's seed.} A person may type a trend: a hashtag, a few
words, a news link or a TikTok link. Hashtags become words to search for
in the news (\texttt{\#FarmearAura} becomes ``Farmear Aura''; reach tags
such as \texttt{\#fyp} are dropped). A news link is the trend's first
article. A TikTok link becomes the trend's video. Its caption is read once
through TikTok's oEmbed, which is public and needs no key, and the video
page itself is never fetched.

\paragraph{Free, and TikTok not scraped.} Every piece of the trends work
is free by rule: public feeds read without a key, the project's own
search engine, and the local models. TikTok is not a signal. On
26~October 2026 its advertisers' trend listing redirected to a page built
in the browser, whose data endpoint answered ``no permission'' without a
signed-in session. Its terms forbid scraping. Its research interface and
the European Union's Digital Services Act Article~40 portal for vetted
researchers are applications measured in months; one published request
for TikTok data through the portal was refused in February 2026. Paid
scrapers would cost money and break the same terms. TikTok therefore
reaches the system as its echo --- the searches and the news a trend
causes --- and as what a person types. Reddit and oEmbed both limit
readers without a key, and on the same day a second request to each came
back empty or refused. The readers treat that as a signal with no trends,
said in the round.

\paragraph{What the signals gave.} On the feeds saved that morning,
Spain's ten trends were a TV talent show, a bullfighter's party, two
tennis matches and a basketball player, a lorry crash, warships off
Ceuta, a television presenter, Madrid's Metro line~10 reopening and a
Netflix adaptation of \emph{La Regenta}: none good news in itself.
r/UpliftingNews linked 24 articles from 20 domains, about 20 of them
positive. A live round over every signal that afternoon read 53 trends
and 92 articles and offered 39 candidates in 13.4~s; the measurement of
what those candidates are worth is under RQ1
(Section~\ref{sec:results-rq1}).
